% !PRINTED 328
\weberpage{328}
%<<CONT
between $0$ and $\xi'$, and another part outside this interval.

From this it now follows generally, by applying the linear transformation, \gesperrt{that on the circle carrying the points $\alpha, x, \xi, \xi'$, if $\xi, \xi'$ are connected by (16), at least one of the roots $\alpha$ lies in each of the two parts into which the circle is divided by $\xi$ and $\xi'$.} The point $x$ and one of the points $\xi, \xi'$ are arbitrary; the third point is determined by (16).

We shall in particular make the assumption that the point $\xi'$ coincides with one of the two points $X$ determined by the equation (12), and thus set, say,
\weq{17}{
X f'(x) + Y f'(y) + (Xy - Yx)\sqrt{(1-n)H} = 0
}

If we introduce this into (16), by setting $\xi', \eta' = X, Y$ and then removing the factor $Xy - Yx$, we obtain
\weq{18}{
(n-1)\,[\xi f'(x) + \eta f'(y)] + (\xi y - \eta x)\sqrt{(1-n)H} = 0,
}
by which the point $\xi$ is uniquely determined.

If we pass to the inhomogeneous form, we obtain [\S.\,106, (7) and (9)] from (17) and (18)
\weq{19}{
x - X = \dfrac{n f(x)}{f'(x) + (n-1)\sqrt{\displaystyle f'(x)^2 - \frac{n}{n-1} f(x) f''(x)}},
}
\weq{20}{
x - \xi = \dfrac{n f(x)}{f'(x) + \sqrt{\displaystyle f'(x)^2 - \frac{n}{n-1} f(x) f''(x)}},
}
in which the square root can have either sign.

The position of the points can then be characterised as follows. There is at least one root $\alpha$ such that
\weq{21}{
\xi,\ \alpha,\ X,\ x,
}
follow one another in order of magnitude, and such that no root lies between $x$ and $X$. We can also choose $\alpha$ so that no further root lies between $\alpha$ and $X$. Between $\alpha$ and $\xi$, however, other roots may possibly lie.
